\documentclass[11pt]{article}

% --- Packages ---
\usepackage[a4paper, margin=1in]{geometry}
\usepackage{amsmath, amssymb, amsthm}
\usepackage{mathtools}
\usepackage{enumitem}
\usepackage[
  pdftitle={Mathematical Foundations of Loss Functions in Machine Learning},
  pdfauthor={Likerener},
  pdfsubject={Mathematical notes on loss functions in machine learning},
  pdfkeywords={loss functions, machine learning, empirical risk, Bayes risk, regularization}
]{hyperref}
\usepackage{xcolor}
\usepackage{parskip}
\usepackage[most]{tcolorbox}
\usepackage{comment}

% --- Theorem-like environments ---
\theoremstyle{definition}
\newtheorem{definition}{Definition}
\newtheorem{proposition}[definition]{Proposition}

% --- Colored boxes ---
\newtcolorbox{mathnote}{
  colback=blue!4!white, colframe=blue!40!black,
  title=In Mathematics, fonttitle=\bfseries, sharp corners, boxrule=0.4pt
}
\newtcolorbox{mlnote}{
  colback=orange!5!white, colframe=orange!50!black,
  title=In Machine Learning, fonttitle=\bfseries, sharp corners, boxrule=0.4pt
}

% --- Title ---
\title{Mathematical Foundations of Loss Functions\\ in Machine Learning}
\author{Likerener}
\date{\today}

\usepackage{fancyhdr}
\pagestyle{fancy}
\fancyhf{}
\fancyfoot[L]{\textcopyright\ 2026 Likerener}
\fancyfoot[C]{Mathematical Foundations of Loss Functions in ML}
\fancyfoot[R]{\thepage}
\renewcommand{\headrulewidth}{0pt}

\begin{document}
\maketitle

\section*{Preface}

This document develops the standard mathematical scaffolding around
\emph{loss functions} in machine learning. For each notion -- predictors,
loss, risk, empirical risk, Bayes optimality, regularisation -- we (i)~give
a textbook definition, (ii)~state the principal mathematical properties, and
(iii)~interpret the object in machine-learning practice. 

%%% The document is self-contained: a topic-indexed list of textbook references is provided at the end.

\section{Predictors, Losses, and Risk}

A supervised-learning problem is specified by a data-generating distribution
$\mathcal{P}$ on $\mathcal{X} \times \mathcal{Y}$, a predictor
$f : \mathcal{X} \to \widehat{\mathcal{Y}}$, and a \emph{loss function}
\[
  L : \mathcal{Y} \times \widehat{\mathcal{Y}} \;\longrightarrow\; \mathbb{R}_{\geq 0}.
\]
The quantity to be minimised is the \emph{risk}, i.e.\ the expected loss
under the data distribution:
\[
  R(f) \;:=\; \mathbb{E}_{(X, Y) \sim \mathcal{P}}\!\left[L\bigl(Y, f(X)\bigr)\right]
  \;=\; \int_{\mathcal{X} \times \mathcal{Y}} L\bigl(y, f(x)\bigr)\, \mathrm{d}\mathcal{P}(x, y).
\]

\bigskip
\hrule
\bigskip

\noindent\textbf{Mathematical background.}

\paragraph{Infimum and supremum.}

\begin{definition}[Infimum and supremum]
For a nonempty set $S \subseteq \mathbb{R}$, the \emph{supremum}
$\sup S \in \mathbb{R} \cup \{+\infty\}$ is the least upper bound of $S$:
(i)~$s \leq \sup S$ for every $s \in S$, and (ii)~for any $u < \sup S$
there exists $s \in S$ with $s > u$. The \emph{infimum} is the greatest
lower bound, defined dually as $\inf S := -\sup(-S)$. The completeness
axiom of $\mathbb{R}$ guarantees that $\sup S$ and $\inf S$ exist (possibly
$\pm\infty$). Unlike $\min$ and $\max$, an infimum or supremum need not
itself belong to $S$.
\end{definition}

\paragraph{Measurable spaces and measurable functions.}

\begin{definition}[$\sigma$-algebra, measurable space, measurable function]
A \emph{$\sigma$-algebra} on a set $\Omega$ is a collection
$\mathcal{F} \subseteq 2^{\Omega}$ containing $\Omega$ and closed under
complementation and countable unions. The pair $(\Omega, \mathcal{F})$ is a
\emph{measurable space}. A function
$f : (\Omega_{1}, \mathcal{F}_{1}) \to (\Omega_{2}, \mathcal{F}_{2})$ is
\emph{measurable} if $f^{-1}(A) \in \mathcal{F}_{1}$ for every
$A \in \mathcal{F}_{2}$. For countable spaces we tacitly take
$\mathcal{F} = 2^{\Omega}$ (all subsets are measurable), so measurability
becomes a vacuous condition.
\end{definition}

\paragraph{Probability spaces (Kolmogorov axioms).}

\begin{definition}[Probability space]
A \emph{probability space} is a triple $(\Omega, \mathcal{F}, \mathbb{P})$
consisting of a sample space $\Omega$, a $\sigma$-algebra $\mathcal{F}$ of
events on $\Omega$, and a \emph{probability measure}
$\mathbb{P} : \mathcal{F} \to [0, 1]$ satisfying Kolmogorov's three axioms:
\begin{enumerate}[label=(\roman*), leftmargin=2em]
  \item (\emph{non-negativity}) $\mathbb{P}(A) \geq 0$ for every
        $A \in \mathcal{F}$;
  \item (\emph{normalisation}) $\mathbb{P}(\Omega) = 1$;
  \item (\emph{countable additivity}) for any pairwise disjoint family
        $\{A_{i}\}_{i \in \mathbb{N}} \subseteq \mathcal{F}$,
        \[
          \mathbb{P}\!\biggl(\,\bigcup_{i \in \mathbb{N}} A_{i}\biggr)
          \;=\; \sum_{i \in \mathbb{N}} \mathbb{P}(A_{i}).
        \]
\end{enumerate}
\end{definition}

\paragraph{Random variables and distributions.}

\begin{definition}[Random variable; pushforward distribution; joint and marginal]
A \emph{random variable} taking values in a measurable space
$(\mathcal{X}, \mathcal{F}_{\mathcal{X}})$ is a measurable function
$X : (\Omega, \mathcal{F}) \to (\mathcal{X}, \mathcal{F}_{\mathcal{X}})$. Its
\emph{distribution} (or \emph{law}) is the \emph{pushforward measure}
$\mathcal{P}_{X}$ on $\mathcal{X}$ defined by
\[
  \mathcal{P}_{X}(A) \;:=\; \mathbb{P}\bigl(X^{-1}(A)\bigr)
  \;=\; \mathbb{P}(X \in A), \qquad A \in \mathcal{F}_{\mathcal{X}}.
\]
For two random variables $X, Y$ on the same probability space, their
\emph{joint distribution} $\mathcal{P}_{X, Y}$ is the pushforward of
$\mathbb{P}$ by the map $\omega \mapsto (X(\omega), Y(\omega))$. The
\emph{marginal distribution} of $X$ is recovered by integrating out $Y$:
$\mathcal{P}_{X}(A) = \mathcal{P}_{X, Y}(A \times \mathcal{Y})$.
\end{definition}

\paragraph{Independence.}

\begin{definition}[Independent random variables]
Two random variables $X, Y$ are \emph{independent} if their joint
distribution factorises into the product of marginals:
\[
  \mathcal{P}_{X, Y}(A \times B) \;=\; \mathcal{P}_{X}(A)\,\mathcal{P}_{Y}(B)
\]
for all measurable $A \subseteq \mathcal{X}$ and $B \subseteq \mathcal{Y}$.
Equivalently, for any pair of events $\{X \in A\}, \{Y \in B\}$,
$\mathbb{P}(X \in A,\, Y \in B) = \mathbb{P}(X \in A)\,\mathbb{P}(Y \in B)$.
A family $X_{1}, \dots, X_{n}$ is \emph{mutually independent} if the joint
distribution factorises into the $n$-fold product of marginals.
\end{definition}

\paragraph{Lebesgue integral.}

\begin{definition}[Lebesgue integral]
The \emph{Lebesgue integral} of a measurable function
$f : (\Omega, \mathcal{F}) \to \mathbb{R}_{\geq 0}$ with respect to
$\mathbb{P}$ is constructed in three stages.
\begin{enumerate}[label=(\arabic*), leftmargin=2em]
  \item (\emph{Simple functions}) For
        $f = \sum_{i=1}^{n} a_{i}\, \mathbf{1}_{A_{i}}$ with
        $a_{i} \geq 0$ and disjoint $A_{i} \in \mathcal{F}$,
        $\displaystyle \int_{\Omega} f\, \mathrm{d}\mathbb{P}
         \;:=\; \sum_{i=1}^{n} a_{i}\, \mathbb{P}(A_{i}).$
  \item (\emph{Non-negative functions}) For general non-negative measurable
        $f$,
        \[
          \int_{\Omega} f\, \mathrm{d}\mathbb{P}
          \;:=\; \sup\!\Bigl\{\,
            \textstyle\int_{\Omega} s\, \mathrm{d}\mathbb{P}
            \;:\; s \text{ simple},\ 0 \leq s \leq f \,\Bigr\}.
        \]
  \item (\emph{Signed functions}) For a general measurable $f$, write
        $f = f^{+} - f^{-}$ with $f^{\pm} \geq 0$ and define
        $\int f\, \mathrm{d}\mathbb{P}
         := \int f^{+}\, \mathrm{d}\mathbb{P} - \int f^{-}\, \mathrm{d}\mathbb{P}$,
        provided not both pieces are infinite.
\end{enumerate}
The Lebesgue integral coincides with the Riemann integral whenever both
exist, but is defined for a vastly larger class of functions and satisfies
strong convergence theorems (monotone, dominated, Fatou).
\end{definition}

\paragraph{Expectation and conditional expectation.}

\begin{definition}[Expectation and conditional expectation]
For a real-valued random variable $Z$ on a probability space
$(\Omega, \mathcal{F}, \mathbb{P})$, the \emph{expectation} is the Lebesgue
integral of $Z$ against $\mathbb{P}$:
\[
  \mathbb{E}[Z] \;:=\; \int_{\Omega} Z\, \mathrm{d}\mathbb{P}.
\]
On a countable set with PMF $p$, this reduces to
$\mathbb{E}[Z] = \sum_{z} z\, p(z)$ when the series converges absolutely.
The \emph{conditional expectation} of $Y$ given $X = x$ is
\[
  \mathbb{E}[Y \mid X = x] \;:=\; \int y\, \mathrm{d}\mathcal{P}_{Y \mid X = x}(y),
\]
or $\sum_{y} y\, p(y \mid x)$ in the discrete case. The
\emph{law of total expectation} (tower property) states
\[
  \mathbb{E}[Y] \;=\; \mathbb{E}_{X}\!\bigl[\mathbb{E}[Y \mid X]\bigr].
\]
\end{definition}

\paragraph{Variance.}

\begin{definition}[Variance and standard deviation]
For a real-valued random variable $Z$ with $\mathbb{E}[Z^{2}] < \infty$ and
mean $\mu := \mathbb{E}[Z]$, the \emph{variance} is
\[
  \operatorname{Var}(Z) \;:=\; \mathbb{E}\!\bigl[(Z - \mu)^{2}\bigr]
  \;=\; \mathbb{E}[Z^{2}] - \mu^{2} \;\geq\; 0.
\]
It satisfies $\operatorname{Var}(aZ + b) = a^{2}\operatorname{Var}(Z)$ and,
for independent $Z_{1}, Z_{2}$,
$\operatorname{Var}(Z_{1} + Z_{2}) = \operatorname{Var}(Z_{1}) + \operatorname{Var}(Z_{2})$.
The \emph{standard deviation} is $\sigma(Z) := \sqrt{\operatorname{Var}(Z)}$.
\end{definition}

\paragraph{Random data and predictors.}

\begin{definition}[Data-generating distribution and predictor]
Let $(\mathcal{X}, \mathcal{F}_{\mathcal{X}})$ and
$(\mathcal{Y}, \mathcal{F}_{\mathcal{Y}})$ be measurable spaces. The
\emph{data-generating distribution} is a probability measure $\mathcal{P}$ on
the product space $\mathcal{X} \times \mathcal{Y}$; we write
$(X, Y) \sim \mathcal{P}$. A \emph{predictor} (or \emph{hypothesis}) is a
measurable function $f : \mathcal{X} \to \widehat{\mathcal{Y}}$, where the
\emph{prediction space} $\widehat{\mathcal{Y}}$ may differ from $\mathcal{Y}$
(e.g.\ probabilistic predictions for classification take values in the
simplex $\Delta(\mathcal{Y})$). The collection of admissible predictors is
the \emph{hypothesis class} $\mathcal{H} \subseteq \widehat{\mathcal{Y}}^{\mathcal{X}}$.
\end{definition}

\begin{mathnote}
A predictor is a measurable function, not just a function: measurability
ensures that $f(X)$ is a random variable and that expectations of $L(Y, f(X))$
are well-defined. For finite or countable $\mathcal{X}, \mathcal{Y}$ this is
automatic; for continuous spaces it is a non-trivial regularity assumption.
\end{mathnote}

\begin{mlnote}
In supervised learning, the hypothesis class $\mathcal{H}$ is realised by a
parameterised model family $\{f_\theta : \theta \in \Theta\}$ -- linear
models, decision trees, neural networks. Choosing $\mathcal{H}$ is the
first design decision in any ML problem and trades off expressiveness
against estimability.
\end{mlnote}

\paragraph{Loss functions.}

\begin{definition}[Loss function]
A \emph{loss function} on target space $\mathcal{Y}$ and prediction space
$\widehat{\mathcal{Y}}$ is a measurable map
\[
  L : \mathcal{Y} \times \widehat{\mathcal{Y}} \;\to\; \mathbb{R}_{\geq 0},
\]
quantifying the cost incurred by predicting $\hat y$ when the true target
is $y$. By convention $L(y, \hat y) = 0$ when the prediction exactly matches
the target.
\end{definition}

\begin{mathnote}
The codomain is taken non-negative for two reasons: (i)~it makes the risk
$\mathbb{E}[L]$ well-defined as a Lebesgue integral, and (ii)~it ensures
that the optimisation problem $\min_f \mathbb{E}[L(Y, f(X))]$ has a
well-defined infimum. Some authors allow signed losses, but then one must
work with the integral of the positive and negative parts separately.
\end{mathnote}

\begin{mlnote}
A loss is the operational specification of "what we are trying to do." A
neural network with a different loss is a different algorithm -- the same
architecture trained with cross-entropy vs.\ squared loss behaves
fundamentally differently. Choosing the loss is therefore the second design
decision, immediately after $\mathcal{H}$.
\end{mlnote}

\paragraph{Risk.}

\begin{definition}[Risk / generalisation error]
The \emph{risk} (or \emph{expected loss}, or \emph{generalisation error}) of
a predictor $f$ under the data-generating distribution $\mathcal{P}$ is
\[
  R(f) \;:=\; \mathbb{E}_{(X, Y) \sim \mathcal{P}}\!\left[L\bigl(Y, f(X)\bigr)\right]
  \;=\; \int_{\mathcal{X} \times \mathcal{Y}} L\bigl(y, f(x)\bigr) \,\mathrm{d}\mathcal{P}(x, y).
\]
\end{definition}

\begin{mathnote}
$R(\cdot)$ is a functional on the hypothesis class, $R : \mathcal{H} \to
\mathbb{R}_{\geq 0}$. By the law of total expectation it factorises as
\[
  R(f) \;=\; \mathbb{E}_{X}\!\left[
    \mathbb{E}_{Y \mid X}\!\left[L(Y, f(X)) \,\big|\, X\right]
  \right],
\]
which separates the dependence on the marginal $\mathcal{P}_X$ from the
conditional $\mathcal{P}_{Y \mid X}$. This decomposition is the
starting-point for the Bayes-optimal analysis in Section~6.
\end{mathnote}

\begin{mlnote}
$R(f)$ is what we ultimately care about -- the average loss on \emph{future}
data drawn from the same distribution. It is unobservable because
$\mathcal{P}$ is unknown; everything in supervised learning is an attempt
to estimate or upper-bound this quantity from finite data.
\end{mlnote}

\section{Empirical Risk and Empirical Risk Minimisation}

Since $\mathcal{P}$ is unknown, $R(f)$ cannot be evaluated directly. Given an
i.i.d.\ sample $\mathcal{D}_n = \{(x_i, y_i)\}_{i=1}^{n}$ from $\mathcal{P}$,
the \emph{empirical risk} is the sample-mean estimator of $R(f)$:
\[
  \widehat{R}_n(f) \;:=\; \frac{1}{n} \sum_{i=1}^{n} L\bigl(y_i, f(x_i)\bigr).
\]
\emph{Empirical Risk Minimisation} (ERM) selects the predictor minimising
this empirical surrogate over a hypothesis class $\mathcal{H}$:
\[
  \widehat{f}_n \;\in\; \arg\min_{f \in \mathcal{H}} \widehat{R}_n(f).
\]
The \emph{generalisation gap} is the discrepancy $R(f) - \widehat{R}_n(f)$
between true and empirical risk.

\bigskip
\hrule
\bigskip

\noindent\textbf{Mathematical background.}

\paragraph{Independent and identically distributed samples.}

\begin{definition}[i.i.d.\ random variables]
Random variables $Z_{1}, \dots, Z_{n}$ are \emph{independent and identically
distributed} (i.i.d.) under a distribution $\mathcal{P}$ if
\begin{enumerate}[label=(\roman*), leftmargin=2em]
  \item each $Z_{i}$ has marginal distribution $\mathcal{P}$;
  \item the joint distribution factorises:
        $\mathbb{P}(Z_{1} \in A_{1}, \dots, Z_{n} \in A_{n})
        = \prod_{i=1}^{n} \mathbb{P}(Z_{i} \in A_{i})$ for all measurable
        $A_{1}, \dots, A_{n}$.
\end{enumerate}
The joint distribution of $(Z_{1}, \dots, Z_{n})$ is the
\emph{$n$-fold product measure} $\mathcal{P}^{\otimes n}$.
\end{definition}

\paragraph{Empirical risk.}

\begin{definition}[Empirical risk]
Given an i.i.d.\ training sample
$\mathcal{D}_n = \{(x_i, y_i)\}_{i=1}^{n} \sim \mathcal{P}^{\otimes n}$, the
\emph{empirical risk} of a predictor $f$ is
\[
  \widehat{R}_n(f) \;:=\; \frac{1}{n} \sum_{i=1}^{n} L\bigl(y_i, f(x_i)\bigr).
\]
\end{definition}

\begin{proposition}[Unbiasedness and variance of $\widehat{R}_n$]
For any fixed $f$,
\[
  \mathbb{E}\bigl[\widehat{R}_n(f)\bigr] = R(f),
  \qquad
  \operatorname{Var}\bigl(\widehat{R}_n(f)\bigr) = \frac{\operatorname{Var}(L(Y, f(X)))}{n}.
\]
\end{proposition}

\paragraph{Modes of stochastic convergence.}

\begin{definition}[Almost sure, in probability, in distribution]
Let $Z_{1}, Z_{2}, \dots$ be real-valued random variables and $Z$ another
such. The sequence \emph{converges to $Z$}:
\begin{itemize}[leftmargin=2em]
  \item \emph{almost surely}, written $Z_n \xrightarrow{\text{a.s.}} Z$, if
        $\mathbb{P}\bigl(\lim_{n \to \infty} Z_n = Z\bigr) = 1$;
  \item \emph{in probability}, written $Z_n \xrightarrow{\;p\;} Z$, if for
        every $\varepsilon > 0$,
        $\mathbb{P}(|Z_n - Z| > \varepsilon) \to 0$ as $n \to \infty$;
  \item \emph{in distribution}, written $Z_n \xrightarrow{\;d\;} Z$, if
        $\mathbb{E}[g(Z_n)] \to \mathbb{E}[g(Z)]$ for every bounded
        continuous function $g$.
\end{itemize}
The implications are
$\xrightarrow{\text{a.s.}} \;\Rightarrow\; \xrightarrow{\;p\;} \;\Rightarrow\; \xrightarrow{\;d\;}$;
the reverse implications fail in general.
\end{definition}

\begin{proposition}[Strong and weak laws of large numbers]
Let $Z_1, Z_2, \dots$ be i.i.d.\ with $\mathbb{E}[|Z_1|] < \infty$ and mean
$\mu = \mathbb{E}[Z_1]$. Then the sample mean
$\bar{Z}_n := \tfrac{1}{n}\sum_{i=1}^n Z_i$ satisfies
$\bar{Z}_n \xrightarrow{\text{a.s.}} \mu$ (\emph{strong LLN}), which implies
$\bar{Z}_n \xrightarrow{\;p\;} \mu$ (\emph{weak LLN}).
\end{proposition}

\paragraph{Hoeffding's inequality.}

\begin{proposition}[Hoeffding's inequality]
Let $Z_{1}, \dots, Z_{n}$ be independent random variables with
$Z_{i} \in [a_{i}, b_{i}]$ almost surely, and let
$\bar{Z}_n = \tfrac{1}{n}\sum_{i=1}^n Z_i$. Then for any $t > 0$,
\[
  \mathbb{P}\!\Bigl(\,\bigl|\bar{Z}_n - \mathbb{E}[\bar{Z}_n]\bigr| > t \,\Bigr)
  \;\leq\; 2 \exp\!\Biggl(-\frac{2 n^{2} t^{2}}{\sum_{i=1}^n (b_i - a_i)^{2}}\Biggr).
\]
Applied to $Z_i = L(y_i, f(x_i))$ with all $Z_i \in [0, B]$, this gives the
empirical-risk concentration bound
\[
  \mathbb{P}\!\Bigl(\,\bigl|\widehat{R}_n(f) - R(f)\bigr| > t \,\Bigr)
  \;\leq\; 2 \exp\!\left(-\frac{2 n t^{2}}{B^{2}}\right).
\]
\end{proposition}

\begin{mathnote}
$\widehat{R}_n$ is the sample-mean estimator of $R$, so by the law of large
numbers $\widehat{R}_n(f) \to R(f)$ almost surely as $n \to \infty$. The
$\mathcal{O}(1/\sqrt{n})$ concentration rate is the foundation of every
finite-sample generalisation bound in statistical learning theory.
\end{mathnote}

\begin{mlnote}
The empirical risk is what your training loop actually reports -- it is the
average per-example loss on the training set. Mini-batch SGD adds a second
stochastic approximation: at each step, $\widehat{R}_n$ is itself replaced
by an average over a mini-batch, yielding an unbiased gradient estimate of
$\widehat{R}_n$.
\end{mlnote}

\paragraph{The $\arg\min$ and $\arg\max$ operators.}

\begin{definition}[$\arg\min$ and $\arg\max$]
For a function $g : X \to \mathbb{R}$ on a nonempty set $X$, the
\emph{arg-minimum} and \emph{arg-maximum} are the (possibly empty, possibly
multi-valued) sets
\[
  \arg\min_{x \in X} g(x) \;:=\; \bigl\{ x_{\ast} \in X \;:\;
    g(x_{\ast}) \leq g(x)\ \text{for all } x \in X \bigr\},
\]
\[
  \arg\max_{x \in X} g(x) \;:=\; \bigl\{ x_{\ast} \in X \;:\;
    g(x_{\ast}) \geq g(x)\ \text{for all } x \in X \bigr\}.
\]
They are related by
$\arg\min_{x} g(x) = \arg\max_{x}\!\bigl(-g(x)\bigr)$, and both exist (are
nonempty) whenever $X$ is finite or whenever $g$ is continuous and $X$ is
compact. When the set is a singleton, we abuse notation and identify
$\arg\min$ (resp.\ $\arg\max$) with its unique element.
\end{definition}

\paragraph{Empirical risk minimisation and the generalisation gap.}

\begin{definition}[ERM]
The \emph{empirical risk minimiser} over a hypothesis class $\mathcal{H}$ is
\[
  \widehat{f}_n \;\in\; \arg\min_{f \in \mathcal{H}} \widehat{R}_n(f).
\]
The \emph{generalisation gap} of a predictor $f$ is
$R(f) - \widehat{R}_n(f)$.
\end{definition}

\begin{mathnote}
Bounding the generalisation gap uniformly over $\mathcal{H}$ -- that is,
controlling $\sup_{f \in \mathcal{H}} |R(f) - \widehat{R}_n(f)|$ -- is the
central problem of statistical learning theory. Capacity measures such as
VC dimension, Rademacher complexity, and PAC-Bayes bounds all serve this
purpose.
\end{mathnote}

\begin{mlnote}
Most training procedures (gradient descent on a loss) are approximations of
ERM. Overfitting is the empirical phenomenon by which
$\widehat{R}_n(\widehat{f}_n)$ becomes small while $R(\widehat{f}_n)$
remains large, i.e.\ the generalisation gap is large. This motivates
regularisation (\S 7).
\end{mlnote}

\section{Regression Losses}

For real-valued targets $\mathcal{Y} = \widehat{\mathcal{Y}} = \mathbb{R}$,
the three canonical losses are
\[
  L_{\mathrm{sq}}(y, \hat y) \;=\; (y - \hat y)^{2}
  \qquad \text{(squared loss)},
\]
\[
  L_{\mathrm{abs}}(y, \hat y) \;=\; |y - \hat y|
  \qquad \text{(absolute loss)},
\]
\[
  L_{\mathrm{Huber}}^{\delta}(y, \hat y) \;=\;
  \begin{cases}
    \tfrac{1}{2}(y - \hat y)^{2}, & |y - \hat y| \leq \delta,\\[2pt]
    \delta\,|y - \hat y| \,-\, \tfrac{1}{2}\delta^{2}, & |y - \hat y| > \delta,
  \end{cases}
  \qquad \text{(Huber loss with threshold } \delta > 0\text{)}.
\]
All three are convex; the first is $C^{\infty}$, the second is non-smooth
at zero, and the third is the unique $C^{1}$ interpolation between them.

\bigskip
\hrule
\bigskip

\noindent\textbf{Mathematical background.}

\paragraph{Convex and strictly convex functions.}

\begin{definition}[Convex, strictly convex, smooth]
A function $g : C \to \mathbb{R}$ on a convex set $C \subseteq \mathbb{R}^{d}$
is \emph{convex} if for all $x, x' \in C$ and $\lambda \in [0, 1]$
\[
  g\bigl(\lambda x + (1 - \lambda) x'\bigr)
  \;\leq\; \lambda\, g(x) + (1 - \lambda)\, g(x').
\]
It is \emph{strictly convex} if the inequality is strict whenever $x \neq x'$
and $\lambda \in (0, 1)$, and is \emph{$L$-smooth} (with constant $L > 0$) if
$\nabla g$ is $L$-Lipschitz: $\|\nabla g(x) - \nabla g(x')\| \leq L \|x - x'\|$.
For twice-differentiable $g$ on a convex domain, convexity is equivalent to 
$\nabla^2 g(x)\succeq 0$. If $\nabla^2 g(x)\succ 0$ everywhere, then $g$ is strictly convex, but the converse need not hold.
\end{definition}

\begin{mathnote}
Convexity is the structural property under which any local minimum of $g$
is also a global minimum, and the gradient $\nabla g$ has the variational
characterisation
$g(x') \geq g(x) + \langle \nabla g(x),\, x' - x\rangle$. These two facts
underlie essentially all of convex optimisation.
\end{mathnote}

\begin{mlnote}
A loss function inherits convexity from being convex in its prediction
argument $\hat y$. Convex losses give well-behaved optimisation -- no spurious
local minima, monotone progress under gradient descent -- but real-world ML
losses (e.g.\ deep-network training losses) are typically non-convex in the
parameters $\theta$ even when the underlying loss $L(y, \cdot)$ is convex.
\end{mlnote}

\paragraph{Smoothness class $C^{k}$.}

\begin{definition}[$C^{k}$ class]
A function $g : U \to \mathbb{R}$ on an open set $U \subseteq \mathbb{R}^{d}$
belongs to the class $C^{k}$, for $k \in \{0, 1, 2, \dots, \infty\}$, if it
has continuous partial derivatives up to order $k$. In particular, $C^{0}$
means continuous, $C^{1}$ means continuously differentiable, $C^{2}$ means
twice continuously differentiable, and $C^{\infty}$ means smooth
(infinitely differentiable). For a single real variable, $g \in C^{k}$
iff $g, g', \dots, g^{(k)}$ all exist and are continuous.
\end{definition}

\paragraph{Squared loss.}

\begin{definition}[Squared loss]
\[
  L_{\mathrm{sq}}(y, \hat y) \;:=\; (y - \hat y)^{2}.
\]
\end{definition}

\begin{mathnote}
$L_{\mathrm{sq}}$ is the squared distance under the Euclidean metric on
$\mathbb{R}$. As a function of $\hat y$ it is strictly convex, infinitely
differentiable, and has constant second derivative $2$; the corresponding
risk $R(f) = \mathbb{E}[(Y - f(X))^2]$ is a quadratic functional. By
Pythagoras of $L^2(\mathcal{P})$, the Bayes-optimal predictor is the
conditional mean $f^\ast(x) = \mathbb{E}[Y \mid X = x]$ (see \S 6).
\end{mathnote}

\begin{mlnote}
Squared loss is the default for regression: it yields closed-form solutions
in linear models (ordinary least squares), Gaussian likelihood
interpretation, and tractable optimisation in neural nets (MSE loss).
Its main weakness is high sensitivity to outliers, since each error is
squared.
\end{mlnote}

\paragraph{Absolute loss.}

\begin{definition}[Absolute loss]
\[
  L_{\mathrm{abs}}(y, \hat y) \;:=\; |y - \hat y|.
\]
\end{definition}

\begin{mathnote}
$L_{\mathrm{abs}}$ is convex but not differentiable at $\hat y = y$. The
Bayes-optimal predictor for absolute loss is the conditional median
$\mathrm{Median}(Y \mid X = x)$. The corresponding empirical procedure is
least-absolute-deviations (LAD) regression, which is more robust to outliers
than OLS but requires linear-programming solvers.
\end{mathnote}

\begin{mlnote}
Absolute loss is used when residuals are heavy-tailed or when one wants
robustness to outliers. Quantile loss
$L_\tau(y, \hat y) := \tau (y - \hat y)_+ + (1 - \tau)(\hat y - y)_+$
generalises it: the Bayes-optimal predictor becomes the $\tau$-th conditional
quantile, enabling quantile regression.
\end{mlnote}

\paragraph{Huber loss.}

\begin{definition}[Huber loss]
For a threshold $\delta > 0$,
\[
  L_{\mathrm{Huber}}^{\delta}(y, \hat y) \;:=\;
  \begin{cases}
    \tfrac{1}{2}(y - \hat y)^{2} & \text{if } |y - \hat y| \leq \delta, \\[2pt]
    \delta\,|y - \hat y| \,-\, \tfrac{1}{2}\delta^{2} & \text{if } |y - \hat y| > \delta.
  \end{cases}
\]
\end{definition}

\begin{mathnote}
Huber loss is the convex combination -- in fact the unique $C^{1}$ function
-- that is quadratic for small residuals and linear for large ones. It is
$C^{1}$ but not $C^{2}$; its derivative is bounded by $\delta$, providing
gradient clipping for free.
\end{mathnote}

\begin{mlnote}
Huber loss interpolates between squared loss (efficient near the optimum)
and absolute loss (robust to outliers). It is the standard choice in
robust regression and is used in DQN-style reinforcement learning to
stabilise Bellman-error targets.
\end{mlnote}

\section{Classification Losses}

For discrete targets, the two canonical losses are
\[
  L_{0\text{-}1}(y, \hat y) \;=\; \mathbf{1}\{\hat y \neq y\}
  \;=\; \begin{cases} 0, & \hat y = y,\\ 1, & \hat y \neq y, \end{cases}
  \qquad \text{(0-1 loss)},
\]
\[
  L_{\mathrm{hinge}}(y, \hat y) \;=\; \max\{0,\; 1 - y\hat y\}
  \qquad \text{(hinge loss, for } \mathcal{Y} = \{-1, +1\},\ \hat y \in \mathbb{R}\text{)}.
\]
$L_{0\text{-}1}$ is the quantity we ultimately wish to minimise (the
misclassification rate), but is non-convex and non-differentiable;
$L_{\mathrm{hinge}}$ is a convex \emph{surrogate} that upper-bounds it under
sign decoding.

\bigskip
\hrule
\bigskip

\noindent\textbf{Mathematical background.}

\paragraph{Indicator function.}

\begin{definition}[Indicator function]
For a set $X$ and a predicate $P(x)$ on $X$ (equivalently, a subset
$A \subseteq X$), the \emph{indicator function} is
\[
  \mathbf{1}\{P(x)\} \;=\; \mathbf{1}_{A}(x) \;:=\;
  \begin{cases} 1, & \text{if } P(x) \text{ is true (equivalently, } x \in A), \\
                0, & \text{otherwise.} \end{cases}
\]
It takes values in $\{0, 1\}$, and its expectation under a distribution $p$
is the probability of the predicate:
$\mathbb{E}_{x \sim p}[\mathbf{1}\{P(x)\}] = \mathbb{P}(P(X))$.
\end{definition}

\paragraph{Subgradient and subdifferential.}

\begin{definition}[Subdifferential]
For a convex function $g : \mathbb{R}^{d} \to \mathbb{R}$ and a point $x_{0}$,
the \emph{subdifferential} of $g$ at $x_{0}$ is the set
\[
  \partial g(x_{0}) \;:=\; \bigl\{\, v \in \mathbb{R}^{d} \;:\;
    g(x) \geq g(x_{0}) + \langle v,\, x - x_{0}\rangle
    \text{ for all } x \in \mathbb{R}^{d} \,\bigr\}.
\]
Elements of $\partial g(x_{0})$ are called \emph{subgradients}. If $g$ is
differentiable at $x_{0}$, then $\partial g(x_{0}) = \{\nabla g(x_{0})\}$;
at non-smooth points (e.g.\ the kink of $|x|$ at $x = 0$), the
subdifferential is a non-singleton interval. A point $x_{0}$ is a global
minimiser of $g$ iff $0 \in \partial g(x_{0})$.
\end{definition}

\paragraph{Zero-one loss.}

\begin{definition}[0-1 loss]
\[
  L_{0\text{-}1}(y, \hat y) \;:=\; \mathbf{1}\{\hat y \neq y\}
  \;=\; \begin{cases} 0 & \hat y = y, \\ 1 & \hat y \neq y. \end{cases}
\]
\end{definition}

\begin{mathnote}
$L_{0\text{-}1}$ is non-convex, non-continuous, and almost-everywhere flat;
its gradient is zero wherever defined. Direct minimisation of empirical 0-1
risk is NP-hard for most non-trivial hypothesis classes. The Bayes-optimal
predictor is the conditional mode $f^{\ast}(x) = \arg\max_{y}
\mathcal{P}(Y = y \mid X = x)$.
\end{mathnote}

\begin{mlnote}
$L_{0\text{-}1}$ is what we ultimately want to minimise for classification --
it is exactly the misclassification rate -- but the gradient-zero property
makes it unusable in gradient-based training. This motivates
\emph{surrogate losses}: convex upper bounds on $L_{0\text{-}1}$ that are
tractable to optimise.
\end{mlnote}

\paragraph{Hinge loss.}

\begin{definition}[Hinge loss]
For $\mathcal{Y} = \{-1, +1\}$ and real-valued predictions
$\hat y \in \mathbb{R}$,
\[
  L_{\mathrm{hinge}}(y, \hat y) \;:=\; \max\{0,\; 1 - y \hat y\}.
\]
\end{definition}

\begin{mathnote}
Hinge loss is convex, piecewise linear, and a surrogate (in fact, an upper
bound) for $L_{0\text{-}1}$ under the sign-decoding rule $\hat{y}_{\text{class}} =
\operatorname{sign}(\hat y)$. Its subdifferential at the kink $y \hat y = 1$
contains $0$, which gives rise to the support-vector phenomenon: only points
near the margin contribute non-zero subgradients.
\end{mathnote}

\begin{mlnote}
Hinge loss is the defining loss of the support-vector machine (SVM). The
"margin" $y \hat y \geq 1$ inside which the loss is zero is the geometric
origin of the SVM's max-margin principle. Multi-class extensions
(Crammer--Singer hinge) generalise this to $\mathcal{Y} = \{1, \dots, K\}$.
\end{mlnote}

\section{Probabilistic Losses: NLL, Cross-Entropy, KL}

When the prediction is a probability distribution
$\hat p \in \Delta(\mathcal{Y})$, three information-theoretic losses arise:
\[
  L_{\mathrm{NLL}}(y, \hat p) \;=\; -\log \hat p(y)
  \qquad \text{(negative log-likelihood)},
\]
\[
  H(p, \hat p) \;=\; -\sum_{y \in \mathcal{Y}} p(y) \log \hat p(y)
  \qquad \text{(cross-entropy)},
\]
\[
  \mathrm{KL}(p \,\|\, \hat p)
  \;=\; \sum_{y \in \mathcal{Y}} p(y) \log \frac{p(y)}{\hat p(y)}
  \qquad \text{(KL divergence)}.
\]
They are linked by the identity
$H(p, \hat p) = H(p) + \mathrm{KL}(p \,\|\, \hat p)$, where $H(p) = -\sum p \log p$
is the entropy. On a one-hot target $p = \delta_{y^{\ast}}$, all three
reduce to $-\log \hat p(y^{\ast})$.

\bigskip
\hrule
\bigskip

\noindent\textbf{Mathematical background.}

\paragraph{Probability simplex.}

\begin{definition}[Probability simplex]
For a countable set $\mathcal{Y}$, the \emph{probability simplex} over
$\mathcal{Y}$ is the set of all probability mass functions on it:
\[
  \Delta(\mathcal{Y}) \;:=\; \biggl\{ p : \mathcal{Y} \to [0, 1] \;\Big|\;
    \sum_{y \in \mathcal{Y}} p(y) = 1 \biggr\}.
\]
For $|\mathcal{Y}| = K$ finite, $\Delta(\mathcal{Y})$ is the $(K-1)$-dimensional
standard simplex embedded in $\mathbb{R}^{K}$:
\(
  \Delta_{K-1}
  = \{p \in \mathbb{R}^{K} : p_k \geq 0,\ \sum_k p_k = 1\}.
\)
\end{definition}

\paragraph{Dirac point mass.}

\begin{definition}[Dirac point mass]
For an element $y^{\ast} \in \mathcal{Y}$, the \emph{Dirac point mass} (or
\emph{one-hot distribution}) at $y^{\ast}$ is the PMF
\[
  \delta_{y^{\ast}}(y) \;:=\; \mathbf{1}\{y = y^{\ast}\}
  \;=\; \begin{cases} 1, & y = y^{\ast},\\ 0, & y \neq y^{\ast}. \end{cases}
\]
It is the extreme point of $\Delta(\mathcal{Y})$ supported at $y^{\ast}$.
On a probability space, the corresponding measure is
$\delta_{y^{\ast}}(A) := \mathbf{1}\{y^{\ast} \in A\}$.
\end{definition}

\paragraph{Gibbs' inequality and Cram\'er--Rao lower bound.}

\begin{proposition}[Gibbs' inequality]
For two PMFs $p, \hat p$ on a countable set $\mathcal{Y}$ with $\hat p(y) > 0$
wherever $p(y) > 0$,
\[
  \mathrm{KL}(p \,\|\, \hat p) \;\geq\; 0,
\]
with equality if and only if $p = \hat p$. Equivalently,
$H(p, \hat p) \geq H(p)$ -- the cross-entropy is bounded below by the entropy
of $p$.
\end{proposition}

\begin{definition}[Fisher information matrix]
For a parametric family $\{p_\theta : \theta \in \Theta \subseteq \mathbb{R}^{d}\}$
satisfying standard regularity conditions (twice differentiability in $\theta$,
support not depending on $\theta$, interchangeability of differentiation and
integration), the \emph{Fisher information matrix} at $\theta$ is
\[
  I(\theta) \;:=\; \mathbb{E}_{Y \sim p_\theta}\!\Bigl[
    \bigl(\nabla_\theta \log p_\theta(Y)\bigr)\bigl(\nabla_\theta \log p_\theta(Y)\bigr)^{\top}
  \Bigr]
  \;=\; -\,\mathbb{E}_{Y \sim p_\theta}\!\bigl[\nabla_\theta^{2} \log p_\theta(Y)\bigr].
\]
\end{definition}

\begin{proposition}[Cram\'er--Rao lower bound]
Under the regularity conditions above, any unbiased estimator
$\widehat{\theta}_n$ of $\theta$ based on $n$ i.i.d.\ observations satisfies
\[
  \operatorname{Cov}(\widehat{\theta}_n) \;\succeq\; \frac{1}{n}\, I(\theta)^{-1},
\]
where $\succeq$ denotes the L\"owner (positive-semidefinite) order. An
unbiased estimator attaining this bound is called \emph{efficient}; the
maximum-likelihood estimator is asymptotically efficient under the same
conditions.
\end{proposition}

\paragraph{Negative log-likelihood.}

\begin{definition}[NLL]
For a target $y \in \mathcal{Y}$ and a predicted distribution
$\hat p \in \Delta(\mathcal{Y})$,
\[
  L_{\mathrm{NLL}}(y, \hat p) \;:=\; -\log \hat p(y).
\]
\end{definition}

\begin{mathnote}
Minimising the empirical NLL over a parametric family $\{\hat p_\theta\}$
is the principle of \emph{maximum likelihood estimation} (MLE):
$\arg\min_\theta \sum_i -\log \hat p_\theta(y_i \mid x_i) = \arg\max_\theta
\sum_i \log \hat p_\theta(y_i \mid x_i)$. Under regularity conditions, the
MLE is consistent, asymptotically normal, and asymptotically efficient
(attains the Cram\'er--Rao lower bound).
\end{mathnote}

\begin{mlnote}
NLL is the most-used loss in modern ML. Language-model pre-training,
classification heads, generative models (VAEs, diffusion, autoregressive
models) all optimise variants of NLL. It is the natural loss whenever the
model outputs a probability distribution.
\end{mlnote}

\paragraph{Cross-entropy.}

\begin{definition}[Cross-entropy]
For two distributions $p, \hat p$ on the same countable set $\mathcal{Y}$,
\[
  H(p, \hat p) \;:=\; -\sum_{y \in \mathcal{Y}} p(y) \log \hat p(y).
\]
\end{definition}

\begin{proposition}[Cross-entropy reduces to NLL on one-hot targets]
If $p = \delta_{y^{\ast}}$ is the point mass at $y^{\ast}$, then
\(
  H(\delta_{y^{\ast}}, \hat p) = -\log \hat p(y^{\ast}) = L_{\mathrm{NLL}}(y^{\ast}, \hat p).
\)
\end{proposition}

\begin{mathnote}
$H(p, \hat p) = H(p) + \mathrm{KL}(p \,\|\, \hat p)$, where $H(p) := -\sum
p \log p$ is the entropy of $p$. Since $H(p)$ does not depend on $\hat p$,
minimising cross-entropy in $\hat p$ is equivalent to minimising the
forward KL divergence $\mathrm{KL}(p \,\|\, \hat p)$.
\end{mathnote}

\begin{mlnote}
In classification, the empirical distribution of labels is one-hot, so
``cross-entropy loss'' and ``NLL'' are used interchangeably in practice.
The PyTorch \texttt{CrossEntropyLoss} module is, technically, the
composition of softmax + NLL with one-hot targets.
\end{mlnote}

\paragraph{Kullback--Leibler divergence as a loss.}

\begin{definition}[KL divergence]
\[
  \mathrm{KL}(p \,\|\, \hat p) \;:=\; \sum_{y} p(y) \log \frac{p(y)}{\hat p(y)}
  \;=\; \mathbb{E}_{y \sim p}\!\left[\log \frac{p(y)}{\hat p(y)}\right].
\]
\end{definition}

\begin{mathnote}
$\mathrm{KL}(p \,\|\, \hat p) \geq 0$ with equality iff $p = \hat p$
(Gibbs' inequality). It is not symmetric:
$\mathrm{KL}(p \,\|\, \hat p) \neq \mathrm{KL}(\hat p \,\|\, p)$ in general,
and the two have qualitatively different behaviours -- ``forward'' KL is
mass-covering, ``reverse'' KL is mode-seeking. The Pythagorean identity
$H(p, \hat p) = H(p) + \mathrm{KL}(p \,\|\, \hat p)$ links cross-entropy
and KL.
\end{mathnote}

\begin{mlnote}
KL appears as a loss in variational inference (ELBO contains
$\mathrm{KL}(q \,\|\, p)$), as a regulariser in RLHF and DPO
(KL-regularised reward maximisation), and as a distillation loss
($\mathrm{KL}(p_{\mathrm{teacher}} \,\|\, p_{\mathrm{student}})$). The
asymmetry matters: choosing forward vs.\ reverse KL changes both the
optimisation landscape and the qualitative shape of the learned distribution.
\end{mlnote}

\section{Bayes Risk and Bayes-Optimal Predictors}

The \emph{Bayes risk} is the smallest achievable risk over all measurable
predictors:
\[
  R^{\ast} \;:=\; \inf_{f \,\text{measurable}} R(f).
\]
A predictor attaining this infimum is called \emph{Bayes-optimal}. For each
canonical loss, the Bayes-optimal predictor is a known functional of the
conditional distribution $\mathcal{P}_{Y \mid X}$:
\[
  L_{\mathrm{sq}} \;\;\Longrightarrow\;\;
  f^{\ast}(x) = \mathbb{E}[Y \mid X = x]
  \qquad \text{(conditional mean)},
\]
\[
  L_{\mathrm{abs}} \;\;\Longrightarrow\;\;
  f^{\ast}(x) = \mathrm{Median}(Y \mid X = x)
  \qquad \text{(conditional median)},
\]
\[
  L_{0\text{-}1} \;\;\Longrightarrow\;\;
  f^{\ast}(x) = \arg\max_{y \in \mathcal{Y}} \mathcal{P}(Y = y \mid X = x)
  \qquad \text{(conditional mode)},
\]
\[
  L_{\mathrm{NLL}} \;\;\Longrightarrow\;\;
  \hat p^{\ast}(\cdot \mid x) = \mathcal{P}(\cdot \mid X = x)
  \qquad \text{(true conditional distribution)}.
\]

\bigskip
\hrule
\bigskip

\noindent\textbf{Mathematical background.}

\begin{definition}[Conditional median and conditional mode]
For a real-valued random variable $Y$ with conditional distribution
$\mathcal{P}_{Y \mid X = x}$, a \emph{conditional median} is any value
$m \in \mathbb{R}$ satisfying
\[
  \mathbb{P}(Y \leq m \mid X = x) \;\geq\; \tfrac{1}{2}
  \qquad \text{and} \qquad
  \mathbb{P}(Y \geq m \mid X = x) \;\geq\; \tfrac{1}{2}.
\]
For discrete $\mathcal{Y}$, the \emph{conditional mode} is any maximiser of
the conditional PMF:
\[
  \mathrm{Mode}(Y \mid X = x) \;\in\; \arg\max_{y \in \mathcal{Y}}
  \mathbb{P}(Y = y \mid X = x).
\]
Neither the conditional median nor the conditional mode is required to be
unique.
\end{definition}

\begin{definition}[Bayes risk and Bayes-optimal predictor]
The \emph{Bayes risk} is the infimum of the risk over all measurable
predictors:
\[
  R^{\ast} \;:=\; \inf_{f \,\text{measurable}} R(f).
\]
A \emph{Bayes-optimal predictor} is any $f^{\ast}$ attaining the infimum,
$R(f^{\ast}) = R^{\ast}$.
\end{definition}

\begin{proposition}[Bayes-optimal predictors for canonical losses]
Under mild regularity,
\begin{enumerate}[label=(\roman*), leftmargin=2em]
  \item \emph{Squared loss:}
        $f^{\ast}(x) = \mathbb{E}[Y \mid X = x]$ (conditional mean).
  \item \emph{Absolute loss:}
        $f^{\ast}(x) = \mathrm{Median}(Y \mid X = x)$ (conditional median).
  \item \emph{0-1 loss:}
        $f^{\ast}(x) = \arg\max_{y \in \mathcal{Y}}
                       \mathcal{P}(Y = y \mid X = x)$ (conditional mode).
  \item \emph{Cross-entropy / NLL:}
        $\hat p^{\ast}(\cdot \mid x) = \mathcal{P}(\cdot \mid X = x)$
        (the true conditional distribution).
\end{enumerate}
\end{proposition}

\begin{mathnote}
Each Bayes-optimal predictor can be derived by minimising the conditional
risk $\mathbb{E}[L(Y, \hat y) \mid X = x]$ pointwise in $\hat y$. For
squared loss this is a quadratic in $\hat y$ minimised at the conditional
mean; for cross-entropy it is a strictly convex functional on the simplex
minimised at the true conditional distribution by Gibbs' inequality.
\end{mathnote}

\begin{mlnote}
Bayes optimality is the irreducible benchmark: no predictor can do better
on the population, no matter how powerful the hypothesis class. The
\emph{excess risk} $R(f) - R^{\ast}$ decomposes into approximation error
($\inf_{\mathcal{H}} R - R^{\ast}$) and estimation error
($R(\widehat{f}_n) - \inf_{\mathcal{H}} R$). Larger $\mathcal{H}$ reduces
the first and increases the second -- the bias--variance trade-off.
\end{mlnote}

\section{Regularised Risk and Structural Risk Minimisation}

Adding a complexity penalty $\Omega(f) \geq 0$ to the empirical risk yields
the \emph{regularised empirical risk}:
\[
  \widehat{R}_{n, \lambda}(f) \;:=\; \widehat{R}_n(f) \;+\; \lambda\, \Omega(f),
  \qquad \lambda \geq 0.
\]
For parametric models $f_\theta$ with $\theta \in \mathbb{R}^{d}$, the
canonical penalties are the $\ell^{p}$ norms,
\[
  \Omega_p(\theta) \;=\; \|\theta\|_p^{p}
  \;=\; \sum_{j=1}^{d} |\theta_j|^{p},
\]
with $p = 2$ (\emph{ridge}, Tikhonov regularisation, weight decay) and
$p = 1$ (\emph{lasso}, sparsity-inducing) by far the most widely used.

\bigskip
\hrule
\bigskip

\noindent\textbf{Mathematical background.}

\paragraph{Regularised empirical risk.}

\begin{definition}[Regularised empirical risk]
For a complexity penalty $\Omega : \mathcal{H} \to \mathbb{R}_{\geq 0}$ and
regularisation strength $\lambda \geq 0$,
\[
  \widehat{R}_{n, \lambda}(f) \;:=\; \widehat{R}_n(f) + \lambda\,\Omega(f).
\]
The \emph{structural risk minimiser} is
$\widehat{f}_{n, \lambda} \in \arg\min_{f \in \mathcal{H}} \widehat{R}_{n, \lambda}(f)$.
\end{definition}

\paragraph{$\ell^{p}$ norms.}

\begin{definition}[$\ell^{p}$ norm]
For $\theta \in \mathbb{R}^{d}$ and $p \in [1, \infty)$, the \emph{$\ell^{p}$
norm} is
\[
  \|\theta\|_{p} \;:=\; \Biggl(\sum_{j=1}^{d} |\theta_j|^{p}\Biggr)^{1/p}.
\]
The limiting cases are
$\|\theta\|_{\infty} := \max_{j} |\theta_{j}|$ and, for $p < 1$,
$\|\cdot\|_{p}$ is a quasi-norm. For $p \geq 1$, $\|\cdot\|_{p}$ satisfies
the three norm axioms (positivity, absolute homogeneity, triangle
inequality), is convex, and is strictly convex iff $p > 1$.
\end{definition}

\paragraph{Probability density functions.}

\begin{definition}[Probability density function]
For a probability measure $\mu$ on $\mathbb{R}^{d}$, a \emph{probability
density function} (PDF) with respect to the Lebesgue measure $\mathrm{d}\theta$
is a non-negative measurable function $p : \mathbb{R}^{d} \to \mathbb{R}_{\geq 0}$
such that
\[
  \mu(A) \;=\; \int_{A} p(\theta)\, \mathrm{d}\theta
  \qquad \text{for every measurable } A \subseteq \mathbb{R}^{d}.
\]
By the Radon--Nikodym theorem, such a density exists if and only if $\mu$
is \emph{absolutely continuous} with respect to the Lebesgue measure -- i.e.\
$\mu(A) = 0$ whenever the Lebesgue measure of $A$ is zero. The total mass
satisfies $\int_{\mathbb{R}^{d}} p(\theta)\, \mathrm{d}\theta = 1$.
\end{definition}

\paragraph{Maximum a posteriori estimation.}

\begin{definition}[MAP estimate]
Given a prior PDF $\pi$ on a parameter $\theta \in \mathbb{R}^{d}$ and a
likelihood $p_\theta(\mathcal{D})$ of observed data $\mathcal{D}$, Bayes'
rule gives the \emph{posterior}
\[
  \pi(\theta \mid \mathcal{D})
  \;=\; \frac{p_\theta(\mathcal{D})\, \pi(\theta)}{\int p_{\theta'}(\mathcal{D})\, \pi(\theta')\, \mathrm{d}\theta'}
  \;\propto\; p_\theta(\mathcal{D})\, \pi(\theta).
\]
The \emph{maximum-a-posteriori (MAP) estimate} is the posterior mode:
\[
  \widehat{\theta}_{\mathrm{MAP}}
  \;:=\; \arg\max_{\theta \in \mathbb{R}^{d}} \pi(\theta \mid \mathcal{D})
  \;=\; \arg\max_{\theta} \Bigl[\,\log p_\theta(\mathcal{D}) \,+\, \log \pi(\theta)\,\Bigr].
\]
When the prior is uniform, $\widehat{\theta}_{\mathrm{MAP}}$ reduces to the
maximum-likelihood estimate.
\end{definition}

\paragraph{Gaussian and Laplace distributions.}

\begin{definition}[Multivariate Gaussian]
The \emph{multivariate Gaussian distribution} $\mathcal{N}(\mu, \Sigma)$ on
$\mathbb{R}^{d}$, with mean $\mu \in \mathbb{R}^{d}$ and positive-definite
covariance $\Sigma \in \mathbb{R}^{d \times d}$, has density
\[
  \mathcal{N}(\theta;\, \mu, \Sigma)
  \;=\; \frac{1}{(2\pi)^{d/2}\, |\Sigma|^{1/2}}\,
  \exp\!\Bigl(-\tfrac{1}{2}(\theta - \mu)^{\top} \Sigma^{-1} (\theta - \mu)\Bigr).
\]
For isotropic covariance $\Sigma = \sigma^{2} I$, the negative log-density
equals $\tfrac{1}{2 \sigma^{2}} \|\theta - \mu\|_{2}^{2}$ up to an additive
constant.
\end{definition}

\begin{definition}[Laplace distribution]
The \emph{Laplace distribution} with location $\mu \in \mathbb{R}$ and
scale $b > 0$ on $\mathbb{R}$ has density
\[
  \mathrm{Laplace}(\theta;\, \mu, b) \;=\; \frac{1}{2 b}
  \exp\!\Bigl(-\frac{|\theta - \mu|}{b}\Bigr).
\]
For an i.i.d.\ Laplace prior on $\theta \in \mathbb{R}^{d}$ with $\mu = 0$,
the negative log-density equals $\tfrac{1}{b}\|\theta\|_{1}$ up to an additive
constant.
\end{definition}

\paragraph{Common penalties (Tikhonov / ridge, lasso).}

\begin{definition}[$\ell^{p}$ penalties]
For a parametric model $f_\theta$ with $\theta \in \mathbb{R}^{d}$, the
$\ell^{p}$ penalty is the $p$-th power of the $\ell^{p}$ norm:
\[
  \Omega_p(\theta) \;:=\; \|\theta\|_{p}^{p}
  \;=\; \sum_{j=1}^{d} |\theta_j|^{p}.
\]
The cases $p = 2$ (\emph{ridge}, Tikhonov) and $p = 1$ (\emph{lasso}) are
canonical.
\end{definition}

\begin{proposition}[Bayesian interpretation]
For a model $\hat p_\theta(y \mid x)$ trained by MLE with NLL, adding a
penalty $\lambda \Omega(\theta)$ corresponds to a maximum a posteriori
(MAP) estimate under a prior $\pi(\theta) \propto \exp(-\lambda \Omega(\theta))$:
\begin{itemize}[leftmargin=2em]
  \item $\Omega(\theta) = \tfrac{1}{2}\|\theta\|_2^{2}$ corresponds to a
        Gaussian prior $\mathcal{N}(0, \lambda^{-1} I)$.
  \item $\Omega(\theta) = \|\theta\|_1$ corresponds to an i.i.d.\ Laplace
        prior with scale $1/\lambda$.
\end{itemize}
\end{proposition}

\begin{mathnote}
$\ell_2$ penalty preserves convexity and differentiability; the regularised quadratic problem
has closed form $\hat{\theta}=(X^\top X+\lambda I)^{-1}X^\top y$ or
$\hat{\theta}=(X^\top X+n\lambda I)^{-1}X^\top y$, depending on the scaling convention.
Under the empirical-risk convention
$\frac1n\sum_i (y_i-x_i^\top\theta)^2+\lambda\|\theta\|_2^2$, the latter is the consistent
form. $\ell_1$ is convex but non-differentiable at the origin; its subdifferential includes
$0$ in a coordinate direction whenever $|\partial \widehat R_n/\partial\theta_j|\le\lambda$,
which produces exact sparsity, i.e. zero coefficients.
\end{mathnote}

\begin{mlnote}
$\ell^2$ regularisation is called \emph{weight decay} in deep learning and
is equivalent (for SGD with fixed learning rate) to multiplicatively
shrinking weights each step. $\ell^1$ regularisation is used for sparse
feature selection. Other "structural" regularisers (dropout, data
augmentation, early stopping) admit similar Bayesian or implicit
interpretations and are studied as instances of structural risk
minimisation.
\end{mlnote}

\bigskip
\hrule
\bigskip

\noindent\textbf{Summary.}
The supervised learning problem has four moving parts:
\begin{itemize}[leftmargin=2em]
  \item a \emph{data distribution} $\mathcal{P}$ on $\mathcal{X} \times \mathcal{Y}$;
  \item a \emph{hypothesis class} $\mathcal{H}$ of measurable predictors;
  \item a \emph{loss function} $L : \mathcal{Y} \times \widehat{\mathcal{Y}}
        \to \mathbb{R}_{\geq 0}$;
  \item a \emph{learning rule} that selects some $\widehat{f} \in \mathcal{H}$
        from data.
\end{itemize}
The loss defines the \emph{risk}
$R(f) = \mathbb{E}_\mathcal{P}[L(Y, f(X))]$, which the learning rule
approximates via the \emph{empirical risk}
$\widehat{R}_n(f) = \tfrac{1}{n}\sum_i L(y_i, f(x_i))$, optionally penalised
by $\lambda \Omega(f)$. The \emph{Bayes risk} $R^{\ast}$ is the
distribution-dependent lower bound on $R$; for canonical losses, the
Bayes-optimal predictor has a closed form -- conditional mean for squared
loss, conditional median for absolute, conditional mode for 0-1, and the
true conditional distribution for cross-entropy.

Every named loss in this document is a different choice within the same
framework. Choosing a loss is, mathematically, choosing which functional of
the conditional distribution $\mathcal{P}_{Y \mid X}$ the Bayes-optimal
predictor recovers; choosing a regulariser is, equivalently, choosing a
prior over the hypothesis class. Every other ML loss (contrastive,
preference-based DPO, RL-style policy-gradient surrogates) can be read as
an instance or extension of this same scaffold.

\section*{Copyright}

\textcopyright\ 2026 Likerener. All rights reserved.

\section*{License}

\textcopyright\ 2026 Likerener. This work is licensed under the Creative Commons Attribution 4.0 International License (CC BY 4.0).

If you use, share, or adapt these notes, please credit the author as follows:

\begin{quote}
Likerener. \textit{Mathematical Foundations of Loss Functions in Machine Learning}. 2026.
\end{quote}

\bigskip
\hrule
\bigskip

\begin{comment}
\section*{References}
\addcontentsline{toc}{section}{References}

The following references are grouped by topic; for any object defined in the
document, one of these texts gives a fuller treatment with proofs.

\paragraph{Measure-theoretic probability.}
P.~Billingsley, \emph{Probability and Measure} (3rd ed., 1995);
R.~Durrett, \emph{Probability: Theory and Examples} (5th ed., 2019);
D.~Williams, \emph{Probability with Martingales} (1991);
O.~Kallenberg, \emph{Foundations of Modern Probability} (2nd ed., 2002).

\paragraph{Real analysis and measure theory.}
W.~Rudin, \emph{Real and Complex Analysis} (3rd ed., 1987);
G.~B.~Folland, \emph{Real Analysis: Modern Techniques and Their Applications}
(2nd ed., 1999).

\paragraph{Statistical inference and estimation theory.}
G.~Casella and R.~L.~Berger, \emph{Statistical Inference} (2nd ed., 2002);
E.~L.~Lehmann and G.~Casella, \emph{Theory of Point Estimation}
(2nd ed., 1998) -- canonical reference for Cram\'er--Rao, Fisher information,
asymptotic efficiency;
L.~Wasserman, \emph{All of Statistics} (2004).

\paragraph{Statistical learning theory.}
L.~Devroye, L.~Gy\"orfi, and G.~Lugosi, \emph{A Probabilistic Theory of
Pattern Recognition} (1996);
S.~Shalev-Shwartz and S.~Ben-David, \emph{Understanding Machine Learning:
From Theory to Algorithms} (2014);
T.~Hastie, R.~Tibshirani, and J.~Friedman, \emph{The Elements of Statistical
Learning} (2nd ed., 2009).

\paragraph{Concentration inequalities.}
S.~Boucheron, G.~Lugosi, and P.~Massart, \emph{Concentration Inequalities:
A Nonasymptotic Theory of Independence} (2013) -- the modern reference for
Hoeffding, Bernstein, McDiarmid, and friends.

\paragraph{Convex analysis and optimisation.}
R.~T.~Rockafellar, \emph{Convex Analysis} (1970) -- canonical reference for
subdifferentials, conjugate duality, and convex functions;
S.~Boyd and L.~Vandenberghe, \emph{Convex Optimization} (2004);
Y.~Nesterov, \emph{Lectures on Convex Optimization} (2nd ed., 2018).

\paragraph{Information theory.}
T.~M.~Cover and J.~A.~Thomas, \emph{Elements of Information Theory}
(2nd ed., 2006) -- canonical reference for entropy, cross-entropy, KL
divergence, and Gibbs' inequality;
D.~J.~C.~MacKay, \emph{Information Theory, Inference, and Learning
Algorithms} (2003).

\paragraph{Bayesian inference and decision theory.}
C.~P.~Robert, \emph{The Bayesian Choice} (2nd ed., 2007);
A.~Gelman, J.~B.~Carlin, H.~S.~Stern, D.~B.~Dunson, A.~Vehtari, and
D.~B.~Rubin, \emph{Bayesian Data Analysis} (3rd ed., 2013);
J.~O.~Berger, \emph{Statistical Decision Theory and Bayesian Analysis}
(2nd ed., 1985) -- for Bayes risk, decision rules, and admissibility.

\paragraph{Probabilistic machine learning.}
C.~M.~Bishop, \emph{Pattern Recognition and Machine Learning} (2006);
K.~P.~Murphy, \emph{Probabilistic Machine Learning: An Introduction} (2022)
and \emph{Probabilistic Machine Learning: Advanced Topics} (2023).

\paragraph{Foundational papers.}
W.~Hoeffding, ``Probability Inequalities for Sums of Bounded Random
Variables,'' \emph{Journal of the American Statistical Association}, 58(301):13--30, 1963;
P.~L.~Bartlett and S.~Mendelson, ``Rademacher and Gaussian Complexities:
Risk Bounds and Structural Results,'' \emph{Journal of Machine Learning
Research}, 3:463--482, 2002.
\end{comment}
\end{document}
